% ============================================================
\section{Implementación}
% ============================================================

\subsection{Caso A: Dependencias, Selección de Dispositivo y Entrenamiento}

La instalación de dependencias en el entorno Python se reduce a dos paquetes
(Roboflow para descargar el dataset y Ultralytics para el modelo):

\begin{lstlisting}[style=code, language=bash]
pip install -q -U roboflow ultralytics
\end{lstlisting}

La selección explícita del dispositivo de cómputo se realiza consultando a
PyTorch; el mismo cuaderno corre en GPU o CPU sin cambios adicionales:

\begin{lstlisting}[style=code, language=Python]
import torch

print("PyTorch:", torch.__version__)
print("CUDA disponible:", torch.cuda.is_available())

device = 0 if torch.cuda.is_available() else "cpu"
\end{lstlisting}

El fine-tuning parte de los pesos preentrenados y apunta al \texttt{data.yaml}
exportado por Roboflow (rutas de \texttt{train/valid/test} y clases):

\begin{lstlisting}[style=code, language=Python]
from ultralytics import YOLO

model = YOLO("yolo26n-seg.pt")

train_results = model.train(
    data=str(YOLO_ROOT / "data.yaml"),
    epochs=60, imgsz=640, batch=16,
    device=device, patience=15,
    project="runs_practica4", name="yolo26n_seg_tools",
    pretrained=True,
)

metrics = model.val(data=str(YOLO_ROOT / "data.yaml"))
\end{lstlisting}

Para el benchmark de inferencia, toda medición sincroniza CUDA antes y después
del cronómetro; omitir esta sincronización subestimaría la latencia real, pues
los lanzamientos de kernels son asíncronos:

\begin{lstlisting}[style=code, language=Python]
def predict_timed(model, images, device, imgsz, half=False):
    if device == "cuda":
        torch.cuda.synchronize()
    t0 = time.perf_counter()
    model.predict(images, device=device, imgsz=imgsz,
                  half=half, verbose=False, conf=0.4)
    if device == "cuda":
        torch.cuda.synchronize()
    return (time.perf_counter() - t0) * 1000.0
\end{lstlisting}

\subsection{Caso B: Scripts de Inferencia en Video}

Cada combinación modelo/dispositivo se ejecuta cambiando dos constantes del
script \texttt{run\_yolo.py}; el FPS instantáneo se calcula por diferencia de
tiempos entre cuadros y se superpone al video anotado:

\begin{lstlisting}[style=code, language=Python]
MODEL_NAME = "yolo12n"   # "yolo12n" o "yolo26n"
DEVICE     = "cpu"       # "cuda" o "cpu"

model = YOLO(str(MODEL_PATH))
results = model.predict(frame, device=DEVICE, conf=0.4, verbose=False)

fps = 1.0 / max(now - prev_t, 1e-6)
\end{lstlisting}

Para Real-ESRGAN (\texttt{run\_superres.py}) el dispositivo se pasa
directamente a PyTorch; en GPU se activa media precisión (FP16) y se procesa
por mosaicos (\texttt{tile=128}) para no exceder la VRAM:

\begin{lstlisting}[style=code, language=Python]
model = RRDBNet(num_in_ch=3, num_out_ch=3, num_feat=64,
                num_block=23, num_grow_ch=32, scale=4)
upsampler = RealESRGANer(
    scale=4, model_path=str(MODEL_PATH), model=model,
    tile=128,
    half=(DEVICE == "cuda"),
    device=torch.device(DEVICE),
)
mejorado, _ = upsampler.enhance(crop, outscale=4)
\end{lstlisting}

\subsection{Caso C: Compilación y Vinculación de Librerías CUDA}

La infraestructura de compilación usa CMake orquestado por un Makefile. En el
laboratorio, OpenCV con CUDA (\texttt{WITH\_CUDA=ON} + \texttt{opencv\_contrib})
está \textbf{instalado a nivel de sistema}, por lo que
\texttt{find\_package(OpenCV)} lo localiza directamente y el objetivo GPU se
habilita de forma automática; los flags de inclusión y enlace contra las
librerías CUDA de OpenCV (\texttt{opencv\_cudaimgproc},
\texttt{opencv\_cudafilters}, \texttt{opencv\_cudaarithm}) los propaga CMake a
través de \texttt{\$\{OpenCV\_LIBS\}}:

\begin{lstlisting}[style=code]
cmake_minimum_required(VERSION 3.16)
project(Practica4Parte1C CXX)
set(CMAKE_CXX_STANDARD 17)

# Laboratorio: OpenCV con CUDA instalado en el sistema;
# find_package lo encuentra sin fijar OpenCV_DIR.
find_package(OpenCV REQUIRED)

add_executable(preprocesamiento_cpu src/pipeline_cpu.cpp)
target_link_libraries(preprocesamiento_cpu PRIVATE ${OpenCV_LIBS})
target_compile_options(preprocesamiento_cpu PRIVATE -O2 -Wall -Wextra)

if(OpenCV_CUDA_VERSION)   # solo si la build tiene WITH_CUDA=ON
    add_executable(preprocesamiento_gpu src/pipeline_gpu.cpp)
    target_link_libraries(preprocesamiento_gpu PRIVATE ${OpenCV_LIBS})
    target_compile_options(preprocesamiento_gpu PRIVATE -O2 -Wall -Wextra)
endif()
\end{lstlisting}

El Makefile expone los objetivos de compilación y ejecución:

\begin{lstlisting}[style=code]
all:
	cmake -B build -DCMAKE_BUILD_TYPE=Release -Wno-dev
	cmake --build build -j$(nproc)

run-cpu: all
	./build/preprocesamiento_cpu

run-gpu: all
	./build/preprocesamiento_gpu

clean:
	rm -rf build
\end{lstlisting}

\begin{lstlisting}[style=code, language=bash]
# Compilar y ejecutar en el laboratorio:
make            # configura y compila ambos ejecutables
make run-cpu    # pipeline sobre cv::Mat
make run-gpu    # pipeline sobre cv::cuda::GpuMat
\end{lstlisting}

El núcleo del programa GPU minimiza las transferencias PCIe: el cuadro sube una
sola vez a VRAM, las seis operaciones se encadenan sobre \texttt{GpuMat} y solo
el resultado final regresa a RAM:

\begin{lstlisting}[style=code, language=C++]
d_frame.upload(frame);                                   // CPU -> GPU (unica subida)
cv::cuda::cvtColor(d_frame, d_gris, cv::COLOR_BGR2GRAY);
gaussianFilter->apply(d_gris, d_suavizado);
erodeFilter->apply(d_suavizado, d_erosionado);
dilateFilter->apply(d_erosionado, d_dilatado);
cannyDetector->detect(d_dilatado, d_bordes);
cv::cuda::equalizeHist(d_dilatado, d_ecualizado);
d_resultado.download(resultado);                         // GPU -> CPU (unica bajada)
\end{lstlisting}

Antes de procesar, el programa valida la presencia de un dispositivo CUDA:

\begin{lstlisting}[style=code, language=C++]
if (cv::cuda::getCudaEnabledDeviceCount() == 0) {
    std::cerr << "No se detecto GPU CUDA-capable via OpenCV.\n";
    return 1;
}
cv::cuda::printShortCudaDeviceInfo(cv::cuda::getDevice());
\end{lstlisting}
